% Fachseminar Report (article-style like legacy)
\documentclass[11pt,a4paper]{article}

% Packages
\usepackage[utf8]{inputenc}
\usepackage[T1]{fontenc}
\usepackage{lmodern}
\usepackage[a4paper,margin=1in]{geometry}
\usepackage{amsmath,amssymb,amsfonts,amsthm,mathtools}
\usepackage{graphicx,subcaption,booktabs}
\usepackage{siunitx}
\usepackage{bm}
\usepackage{microtype}
\usepackage{hyperref}
\usepackage[nameinlink]{cleveref}
\usepackage{enumitem}
\usepackage{algorithm}
\usepackage{algpseudocode}

% Graphics paths (adapted to current repo)
\graphicspath{{./}{./plots/}{./../plots/}{./results/}{./../results/}}

% Optional auto macros if present
\InputIfFileExists{auto_macros.tex}{}{}

% Macros
\newcommand{\R}{\mathbb{R}}
\newcommand{\Om}{\Omega}
\newcommand{\om}{\omega}
\newcommand{\half}{\tfrac{1}{2}}
\newcommand{\ip}[2]{\left\langle #1, #2 \right\rangle}
\newcommand{\ipOm}[2]{\left( #1, #2 \right)_{L^2(\Omega)}}
\newcommand{\ipom}[2]{\left( #1, #2 \right)_{L^2(\omega)}}
\newcommand{\normOm}[1]{\left\| #1 \right\|_{L^2(\Omega)}}
\newcommand{\normom}[1]{\left\| #1 \right\|_{L^2(\omega)}}
\newcommand{\normU}[1]{\left\| #1 \right\|_{U}}
\newcommand{\ipU}[2]{\left( #1, #2 \right)_{U}}
\newcommand{\Chi}{\chi}

% Theorem-like environments (style parity with legacy)
\newtheorem{theorem}{Theorem}
\newtheorem{lemma}{Lemma}
\newtheorem{remark}{Remark}

\title{Report — Comparison of Optimization Methods for a PDE-Constrained Optimal Control Problem}
\author{Fachseminar}
\date{\today}

\begin{document}
\maketitle
\tableofcontents

\section{Continuous Problem}
We consider the PDE-constrained optimal control problem on a bounded Lipschitz domain $\Om = (0,1)^2 \subset \R^2$ with a measurable control region $\om \subset \Om$. The state $y\colon \Om \to \R$ and the control $u\colon \om \to \R$ are coupled through the Poisson equation with homogeneous Dirichlet boundary conditions:
\begin{align}
  -\Delta y &= \Chi_{\om} \, u &&\text{in } \Om, \\
  y &= 0 &&\text{on } \partial\Om.
\end{align}
Given a desired state $y_d\in L^2(\Om)$ and a Tikhonov parameter $\beta>0$, we minimize the quadratic objective
\begin{equation}
  F(y,u) = \tfrac12\, \normOm{y - y_d}^2 + \tfrac{\beta}{2}\, \normom{u}^2.
\end{equation}
We adopt the space setting
\begin{equation}
  y \in H_0^1(\Om), \qquad u \in L^2(\om),
\end{equation}
where the right-hand side $\Chi_{\om} u$ is interpreted as an element of $H^{-1}(\Om)$ via the embedding $L^2(\om)\hookrightarrow H^{-1}(\Om)$.

Well-posedness of the state equation follows from Lax–Milgram: for each $u\in L^2(\om)$ there exists a unique $y\in H_0^1(\Om)$ such that the weak form below holds. Existence and uniqueness of an optimal control follow from strict convexity and coercivity of $F$.

\section{Weak Formulation}
Let $V := H_0^1(\Om)$. The weak form of the state equation reads: find $y\in V$ such that
\begin{equation}
  a(y,v) = \ell(v;u) \qquad \forall v\in V,
\end{equation}
with the bilinear form and linear functional given by
\begin{equation}
  a(y,v) := \int_{\Om} \nabla y\cdot\nabla v\,dx, \qquad
  \ell(v;u) := \int_{\om} u\, v\,dx.
\end{equation}
The control-to-state map $S\colon L^2(\om)\to V$ is defined by $y = S u$ as the unique solution of the above variational problem.

To derive first-order optimality conditions, we introduce the adjoint state $p\in V$ defined by
\begin{equation}
  a(w,p) = \int_{\Om} (y - y_d)\, w\,dx \qquad \forall w\in V.
\end{equation}
The reduced formulation in the control (see \Cref{sec:reduced,sec:grad-hess}) yields the optimality condition in $L^2(\om)$: $\nabla j(u)=0$ with an explicit expression for the gradient via $y$ and $p$.

\section{Finite Element Discretization}
We consider a conforming triangulation $\mathcal{T}_h$ of $\Om$ and the continuous, piecewise linear finite element space
\begin{equation}
  V_h := \{ v_h\in C^0(\overline{\Om}) : v_h|_K \text{ is affine } \forall K\in \mathcal{T}_h,\ v_h|_{\partial\Om}=0 \} \subset V.
\end{equation}
For the control, we use a compatible discrete space $U_h\subset L^2(\om)$. In the simplest case, we take the $V_h$ basis restricted to the control region and do not enforce boundary conditions on $U_h$:
\begin{equation}
  U_h := \operatorname{span}\{ \phi_i|_{\om} : \phi_i \text{ nodal basis of } V_h\},
\end{equation}
which is standard when the control acts as a distributed source over $\om$.

The discrete state $y_h\in V_h$ for a given $u_h\in U_h$ solves
\begin{equation}
  a(y_h, v_h) = \ell(v_h; u_h) \qquad \forall v_h\in V_h.
\end{equation}
Homogeneous Dirichlet conditions are imposed on $V_h$ (e.g., by modifying the stiffness matrix rows/columns associated with boundary nodes), while no essential conditions are applied to $U_h$.

\section{Discrete Operators}\label{sec:operators}
Let $\{\phi_i\}_{i=1}^{n_V}$ be the nodal basis of $V_h$ and $\{\psi_j\}_{j=1}^{n_U}$ a basis of $U_h$. We define the usual FE matrices:
\begin{align}
  A_{ij} &:= \int_{\Om} \nabla\phi_i\cdot\nabla\phi_j\,dx && \text{(stiffness on $\Om$)}, \\
  M_{ij} &:= \int_{\Om} \phi_i\,\phi_j\,dx && \text{(mass on $\Om$)}, \\
  (M_U)_{jk} &:= \int_{\om} \psi_j\,\psi_k\,dx && \text{(mass on $\om$)}, \\
  B_{ij} &:= \int_{\om} \psi_j\,\phi_i\,dx && \text{(control coupling)}.
\end{align}
With these definitions, the discrete state and adjoint equations read
\begin{align}
  A y &= B u, \\
  A p &= M (y - y_d),
\end{align}
where $y,p\in\R^{n_V}$ and $u\in\R^{n_U}$ are the coefficient vectors, and $y_d\in\R^{n_V}$ is the $V_h$ projection of the data. Dirichlet boundary conditions are applied to $A$ (and the corresponding right-hand sides) in the standard way; the matrices $M_U$ and $B$ are assembled using cell/subdomain tags so that integrations over $\om$ are accurate.

\section{Reduced Functional}\label{sec:reduced}
Eliminating the state via $y = A^{-1} B u$ yields the discrete reduced cost
\begin{equation}
  j(u) := \tfrac12 (y - y_d)^\top M (y - y_d) + \tfrac{\beta}{2}\, u^\top M_U u,\qquad y = A^{-1} B u.
\end{equation}
Expanding in $u$ gives the quadratic form
\begin{equation}
  j(u) = \tfrac12\, u^\top\underbrace{\big( B^\top A^{-1} M A^{-1} B + \beta M_U \big)}_{\widetilde H}\,u
         - u^\top B^\top A^{-1} M y_d 
         + \tfrac12\, y_d^\top M y_d,
\end{equation}
which shows $j$ is strictly convex for $\beta>0$. The constant term does not affect optimization.

\section{Gradient and Hessian}\label{sec:grad-hess}
We equip the control space with the $U$-inner product
\begin{equation}
  \ipU{a}{b} := a^\top M_U b, \qquad \normU{a}^2 = \ipU{a}{a}.
\end{equation}
For a given $u$, compute the state $y$ and the adjoint $p$ from
\begin{equation}
  A y = B u, \qquad A p = M (y - y_d).
\end{equation}
Then the Riesz-represented gradient in the $U$-inner product is
\begin{equation}
  \nabla j(u) = M_U^{-1} \big( B^\top p + \beta M_U u \big) = M_U^{-1} B^\top p + \beta u.
\end{equation}
Equivalently, the $U$-duality gradient is $g_U(u) := B^\top p + \beta M_U u$. The (Riesz) Hessian operator is the SPD map
\begin{equation}
  H = M_U^{-1} \big( B^\top A^{-1} M A^{-1} B + \beta M_U \big),
\end{equation}
so that for any direction $d$ the Hessian-vector product can be computed by
\begin{align}
  &\text{solve } A\, y'[d] = B\, d, \\
  &\text{solve } A\, p'[d] = M\, y'[d], \\
  &\text{return } H d = M_U^{-1} \big( B^\top p'[d] + \beta M_U d \big).
\end{align}
This structure underpins fixed-step gradient descent with step $1/L$, Barzilai–Borwein steps using the $U$-metric, and Nesterov acceleration with restarts.

% --- End of initial chapters ---

\end{document}
